% Section 2: Literature
\section{Relation to prior work}\label{sec:lit}

\subsection{The OTS-NIDC lineage}

The method descends from Chu \citep{chu2011}, who combined optimal time-step
selection with non-iterative defect correction to boost finite-difference
schemes for linear and semilinear PDEs to fourth order on three-point stencils,
with extensions to more dimensions and irregular domains. Chu and Lambers
\citep{chu2009} extended OTS to spatially \emph{variable} leading coefficients
via an optimal-grid transformation, with the constant-coefficient step
$\dto = \dx^2/6D$ recovered when the coefficient is constant. Razi, Attar and
Vedula \citep{razi2015} coupled grid adaptation with non-iterative defect
correction; our critique notes (\texttt{literature/Razi2015\_critique.md})
observe that their joint optimization degenerates to the OTS step for linear
coupling operators --- a point we return to when justifying why the coupled
corrections of \S\ref{sec:coupled} need no further optimization.

The present work is, to our knowledge, the first systematic extension of this
lineage to \emph{stochastic} equations --- accomplished not by discretizing a
random field but by moving up one level of abstraction: the moment hierarchy of
the SPDE is itself a family of deterministic reaction--diffusion equations, so
the entire deterministic theory applies unmodified (\S\ref{sec:stochastic}).

\subsection{SPDE numerics: the pathwise mainstream}

The numerical SPDE literature is dominated by pathwise schemes. The textbook
frame is Lord, Powell and Shardlow \citep{lord2014}, building on
Kloeden--Platen \citep{kloeden1992}. Within it:

\begin{itemize}
\item \textbf{Euler--Maruyama and Milstein.} The canonical schemes, weak order 1
and strong order $1/2$ (Milstein: strong order 1 for commutative noise)
\citep{kloeden1992,leonhard2014,hajiali2023}. For SPDEs the spatial
semi-discretization error typically dominates; strong rates are further limited
by the spatial regularity of the noise. Every pathwise method pays a Monte Carlo
factor $O(M^{-1/2})$ to reach moments. That is the competitor baseline for this
paper: identical statistics, obtained $M$ times over.

\item \textbf{The Davie--Gaines barrier.} Davie and Gaines \citep{davie2001}
proved that numerical schemes for parabolic SPDEs driven by space-time white
noise, converging through Wiener increments, cannot exceed strong order $1/2$ in
time. This is the fundamental obstruction to high-order \emph{pathwise} schemes
for \eqref{eq:spde} --- and it is precisely what the moment-level reformulation
sidesteps, by never entering the pathwise problem.

\item \textbf{Weak-rate separation.} Bolin, Kirchner and Kov\'acs
\citep{bolin2018} and Cui--Sun \citep{cui2023} exemplify a robust phenomenon:
weak rates can exceed strong rates for SPDE approximations. Our contribution
pushes this to its structural limit: at the moment level, the problem is
deterministic, so the ``weak rate'' is simply the accuracy of a deterministic
solver --- arbitrarily high, at the price of the closure.

\item \textbf{Structure-preserving alternatives.} Exponential integrators treat
the stiff linear part exactly \citep{lord2010}; modified semi-implicit schemes
evolve each Fourier mode exactly as an Ornstein--Uhlenbeck process
\citep{tambue2016}; splitting schemes preserve state-space structure in chemical
Langevin dynamics \citep{jovanovski2025,ulander2024}. These are the honest
competitors for stiff reaction--diffusion SPDEs. They remain pathwise; the
per-mode-exact philosophy, however, is the same one that makes our decoupled
M/V-block stepping work (\S\ref{sec:coupled}).
\end{itemize}

\subsection{Moment methods and closures}

The moment hierarchy approach has a long history in kinetic theory and
population dynamics; the closure problem is its central obstruction. Gaussian
(Wick) closure is exact for linear dynamics driven by Gaussian noise and
otherwise a weak-noise approximation that fails for bimodal or intermittent
responses --- a point developed carefully in the companion primer
(\texttt{moment\_hierarchy\_primer.tex}) and used as the scope boundary here.
The NIDC framework's position is deliberately narrow: it claims nothing about
closures, only that \emph{when a moment system closes, it can be solved to
arbitrary order on cheap stencils}.

\subsection{Multirate time stepping}

Resolving competing optimal steps (\S\ref{sec:coupled}) is a multirate problem.
The classical theory of Gear--Wells \citep{gear1984} and Fok--Rosales
\citep{fok2010} gives the interpolation-order bookkeeping used in
\S\ref{sec:coupled} (cubic Hermite suffices through $p=2$); Richardson-style
extrapolation \citep{oates2024,
debrabant2011} is the post-processing alternative that buys one order at the
cost of two coordinated solves and shared-noise synchronization, whereas NIDC
buys two orders inside a single solve. Recent multirate partitioned Runge--Kutta
methods \citep{roberts2018} are the modern heavy machinery for the same problem
class.

\subsection{Positioning summary}

\begin{center}
\begin{tabular}{@{}lll@{}}
\toprule
Method & Moment accuracy & Cost to moments \\
\midrule
Euler--Maruyama \citep{lord2014} & weak order 1 & $O(M)$ paths, $O(M^{-1/2})$ MC error \\
Milstein \citep{leonhard2014} & weak order $\ge 1$ & same, more work per path \\
Richardson on FD+EM \citep{oates2024} & weak order 2 & two solves, shared noise \\
Exponential integrators \citep{lord2010} & inherits time integrator & pathwise + MC \\
\textbf{OTS-NIDC (this work)} & \textbf{4th order (moment PDE solve)} & \textbf{one deterministic solve, single pass} \\
\bottomrule
\end{tabular}
\end{center}

The comparison is deliberately asymmetric: the pathwise methods deliver
information (paths) that OTS-NIDC does not attempt. The claim is narrower and
cheaper: \emph{if the deliverable is moments, and the noise model closes the
hierarchy, then the moments are a deterministic PDE problem and should be solved
like one.}
